\section{Automaton LR(0) dan SLR(1)}
\label{sec:ch12-slr}

\subsection{Automaton LR(0)}

Automaton LR(0) adalah graf berarah dengan state berupa himpunan
item dan transisi berupa goto. State awal adalah closure item awal.
State penerima adalah state yang memuat item \([S' \to S \cdot]\).

\begin{figure}[H]
\centering
\begin{tikzpicture}[node distance=8mm and 14mm]
  \node[phase, text width=2.4cm] (i0) {\(I_0\)};
  \node[phase, text width=2.4cm, right=of i0] (i2) {\(I_2\)};
  \node[phase, text width=2.4cm, below=of i2] (i4) {\(I_4\)};
  \node[phase, text width=2.4cm, right=of i2] (i7) {\(I_7\)};
  \node[phase, text width=2.4cm, below=of i7] (i10) {\(I_{10}\)};
  \draw[arrow] (i0) -- node[above, font=\small] {\(E\)} (i2);
  \draw[arrow] (i0) -- node[below, font=\small] {\(T\)} (i4);
  \draw[arrow] (i2) -- node[above, font=\small] {\(+\)} (i7);
  \draw[arrow] (i7) -- node[right, font=\small] {\(T\)} (i10);
\end{tikzpicture}
\caption{Bagian automaton LR(0) untuk grammar ekspresi}
\label{fig:ch12-automaton}
\end{figure}

\subsection{Tabel SLR(1)}

SLR(1) menggunakan automaton LR(0) dan himpunan FOLLOW untuk
memutuskan reduksi:
\begin{enumerate}[label=\arabic*.,leftmargin=*]
  \item jika item \([A \to \alpha \cdot a \beta]\) dan \(a\) terminal,
        isi \(\text{ACTION}[I,a] = \text{shift goto}(I,a)\);
  \item jika item \([A \to \alpha \cdot]\) dan \(A \neq S'\), untuk
        setiap \(a \in \text{FOLLOW}(A)\), isi
        \(\text{ACTION}[I,a] = \text{reduce } A \to \alpha\);
  \item jika item \([S' \to S \cdot]\), isi
        \(\text{ACTION}[I,\$] = \text{accept}\);
  \item untuk nonterminal \(A\), isi
        \(\text{GOTO}[I,A] = \text{goto}(I,A)\).
\end{enumerate}

\begin{table}[H]
\centering
\caption{Potongan tabel SLR(1) untuk grammar ekspresi}
\label{tab:ch12-slr}
\begin{tabular}{@{}c c c c c@{}}
\toprule
\textbf{State} & \code{id} & \code{+} & \code{*} & \code{\$} \\
\midrule
0 & s5 & & & \\
1 & & & & \\
2 & & s7 & & acc \\
3 & & r(\(T \to F\)) & r(\(T \to F\)) & r(\(T \to F\)) \\
4 & & r(\(E \to T\)) & s8 & r(\(E \to T\)) \\
5 & & r(\(F \to \code{id}\)) & r(\(F \to \code{id}\)) &
r(\(F \to \code{id}\)) \\
\bottomrule
\end{tabular}
\end{table}

\subsection{Alat Bantu}

Skrip \file{slr\_table.py} membangun koleksi kanonik dan tabel
SLR(1). Fungsi \code{closure} dan \code{goto} ditunjukkan pada
\lstref{lst:ch12-closure}.

\includecode[
  language=Python,
  firstline=26,
  lastline=54,
  firstnumber=26,
  caption={Fungsi \code{closure} dan \code{goto}},
  label={lst:ch12-closure}
]{../../code/tools/slr_table.py}

Fungsi \code{build\_slr\_table} memakai FOLLOW untuk mengisi aksi
reduce. Potongannya ditunjukkan pada \lstref{lst:ch12-build}.

\includecode[
  language=Python,
  firstline=89,
  lastline=131,
  firstnumber=89,
  caption={Fungsi \code{build\_slr\_table}},
  label={lst:ch12-build}
]{../../code/tools/slr_table.py}

Alat ini juga melaporkan konflik. Untuk grammar ekspresi, hasilnya
adalah 12 state dan tidak ada konflik.
